\chapter{第二十三章：三类既得利益者的转型困境}


\textit{"既得利益者是最难被说服的群体，因为他们失去的最多。"}
——约翰·肯尼迪

\newpage

老王我年轻的时候，特别烦那种"老油条"——什么都知道，什么都不动，领导说什么都点头，但就是不干活。

后来我自己岁数大了，我发现我开始理解他们了。

他们不是懒，他们是在保护自己。他们在这个位置上干了二十年，靠的就是那点经验、那点关系、那点信息差。变革一来，他们发现自己的那点"资本"可能一文不值了。你让他们拥抱变革？那不是与虎谋皮吗？

我现在明白了：既得利益者不是坏人，只是被时代抛下的人。他们需要的不是批评，是理解。

\section{一、谁是既得利益者}


在AI变革中，最容易被忽视的既得利益者，往往是那些在传统组织中拥有权力和地位的人。

他们不是"坏人"，也不一定是在故意阻挠变革。他们只是那些在现有组织结构中投入最多、受益最多、因此对维持现状最有动力的人。

在AI转型中，有三类既得利益者尤其值得关注：他们各自的利益诉求不同，对变革的威胁感知不同，因此需要不同的应对策略。

\section{二、第一类：中层管理者}


\textbf{他们的困境}

中层管理者是AI时代受影响最大的群体之一。

在传统组织中，中层管理者的价值主要体现在两个功能上：信息汇总和指令传达。他们收集一线的零散信息，加工成有意义的报告，传递给高层；他们把高层的战略意图，转化成一线可以执行的具体指示。

AI正在替代这两个功能。

当AI可以实时汇总来自所有一线员工的信息并生成报告时，"信息枢纽"的价值急剧下降。当AI可以把高层的战略意图自动分解成具体的执行任务时，"指令传达"的价值同样在下降。

但这里有一个微妙的悖论：很多中层管理者自己并没有意识到这一点。他们仍然认为自己作为"管理者"的价值是不言而喻的。

就像以前单位的门卫，掌握着进出人员的"信息"，谁来了找谁都得过他这关。后来有了门禁系统，他的"信息枢纽"价值就没了。但他可能还在想：我在这儿站了二十年，谁进谁出我还不知道吗？

\textbf{他们的应对策略}

面对AI变革，中层管理者通常有几种典型的应对策略：

第一种是主动拥抱：认识到AI将替代传统管理职能，主动学习AI技能，转型成为"AI时代的有效管理者"。

第二种是被动等待：继续维持现状，寄希望于变革不会真正发生，或者等到变革真正威胁到自己时再应对。

第三种是阻挠变革：利用自己对组织的影响力，缓慢或阻挠AI转型的推进。

\textbf{如何帮助他们}

对于组织来说，帮助中层管理者成功转型，往往比替换他们更有效——因为他们在组织中积累了大量的本地知识和关系资本。

具体的帮助方式包括：帮助他们理解AI对他们角色的真实影响，给他们提供学习新技能的机会和实践场景，重新定义他们的价值定位，让他们看到在AI时代同样可以贡献价值。

\section{三、第二类：专业职能部门}


\textbf{他们的困境}

专业职能部门（财务，法务、人力资源等）是另一个在AI时代面临挑战的群体。

在传统组织中，这些部门的权威建立在"专业知识"之上。财务部门有财务知识，法务部门有法律知识，人力资源部门有管理知识。这些知识让他们在组织中拥有建议权和否决权。

AI正在侵蚀这种知识权威。

当AI可以快速提供法律研究、财务分析、HR政策建议时，专业职能部门的知识优势在减弱。他们发现自己的建议权被"AI建议"所挑战。

同时，这些部门的另一层价值——流程执行——也在被AI自动化所替代。

\textbf{他们的应对策略}

专业职能部门面对AI变革的典型反应，通常是强调"AI可能出错"、"AI没有考虑合规要求"、"AI忽略了具体情境"。

这些担忧并不全是错误的——它们在某些情况下是真实的。但这些担忧，有时也是保护部门利益的借口。

就像有的财务老法师会说："AI做账？AI懂什么叫合理避税吗？AI知道这笔账该怎么调吗？"——其实有些AI可能真的比他更懂最新的税法和准则。但他不愿意承认，因为这等于承认自己那几十年经验可能不值钱了。

\textbf{如何帮助他们}

帮助专业职能部门转型，需要重新定义他们的价值。

当知识权威不再稀缺时，专业职能部门需要找到新的价值定位：他们可以成为"AI输出的审核者和判断者"，利用自己的专业知识评估AI建议的合理性和适用性；他们可以成为"复杂情境的判断者"，在AI无法处理的模糊情况下提供人类判断。

\section{四、第三类：IT部门}


\textbf{他们的困境}

你可能会惊讶：IT部门不是应该支持AI变革吗？他们是技术的拥有者。

但事实上，IT部门往往也是AI变革的隐性阻力者。

在传统组织中，IT部门的权力建立在"技术门槛"之上。普通人无法自己获取或部署技术，需要通过IT部门。IT部门控制着技术采购、架构设计、系统运维。

当AI变得更加易用，当业务部门可以直接采购AI服务（Shadow IT）时，IT部门的传统权力基础在动摇。

\textbf{他们的应对策略}

IT部门的典型反应包括：强调AI的安全风险、数据隐私风险、系统集成风险，强调任何新技术的引入都必须经过IT部门的审核和批准。

这些担忧有其合理性——AI确实带来了新的安全和隐私挑战。但有时，这些担忧也被用来保护IT部门在技术决策中的主导地位。

\textbf{如何帮助他们}

帮助IT部门转型，需要让他们从"技术守门人"变成"AI赋能者"。

IT部门可以在AI时代发挥新的价值：他们可以成为组织AI能力的架构师，帮助各个业务部门设计合理的AI部署方案；他们可以成为AI安全和隐私的守护者，建立组织的AI治理框架；他们可以成为AI能力的集成者，确保各种AI系统能够协同工作。

\section{五、既得利益者转型的共同要素}


虽然三类既得利益者的具体困境不同，但他们成功转型有一些共同的要素。

\textbf{第一，重新定义价值。}

当原有的价值基础被AI动摇时，既得利益者需要找到新的价值基础。这个新的价值基础，通常是AI难以替代的人类能力：判断力、关系力、复杂情境处理能力。

\textbf{第二，提供学习机会。}

转型需要新的技能。组织需要提供系统的学习机会，让既得利益者能够学习AI时代所需的新能力。

\textbf{第三，给予实践场景。}

光学习是不够的，需要在实践中应用新技能。组织需要为既得利益者提供可以尝试新角色的实践场景。

\textbf{第四，承认历史贡献。}

在帮助既得利益者转型时，承认他们的历史贡献是重要的。如果转型方案让他们感觉自己的过去被否定，他们会产生防御性反应。

\section{六、结语：阻力来自恐惧，也来自被看见的需要}


既得利益者对变革的阻力，往往来自恐惧——对未知的恐惧，对失去权力的恐惧，对无法适应的恐惧。

理解这一点，有助于我们以同理心对待这种阻力，而不是简单地给他们贴上"阻挠者"的标签。

同时，帮助既得利益者转型，也需要认识到他们被看见的需要。当他们在组织中的重要性被AI所超越时，他们需要找到新的、被看见的方式来做贡献。

这种转型的过程，对既得利益者和组织来说都是艰难的。但它是必要的。

肯尼迪说"既得利益者是最难被说服的群体，因为他们失去的最多"——这话是真的。但最难被说服不等于不能被说服，关键是你得让他们看到：转型不是剥夺，而是重新定义他们的价值。

\newpage

\textbf{本章小结：}

1. 三类既得利益者面临不同的困境：中层管理者（信息枢纽价值被替代）、专业职能部门（知识权威被侵蚀）、IT部门（技术守门人角色被挑战）
2. 每类既得利益者有典型的应对策略：中层管理者（主动拥抱/被动等待/阻挠变革）、专业职能部门（强调风险）、IT部门（强调安全和合规）
3. 帮助转型的共同要素：重新定义价值、提供学习机会、给予实践场景、承认历史贡献
4. 阻力来自恐惧（对未知的恐惧），也来自被看见的需要（需要找到新的贡献方式）
